\chapter{TINJAUAN PUSTAKA}
	\section{Landasan Teori}
		\subsection{Verifikasi Dokumen Akademik}
		\setlength{\leftskip}{1.4cm}
		Dokumen akademik dalam penelitian ini dipandang sebagai bentuk khusus dari informasi yang memiliki fungsi tertentu, seperti sebagai bukti, alat pembuktian, maupun sarana pendukung proses pendidikan\footcite{Stolyarov2021}. Pemanfaatan blockchain untuk verifikasi dokumen akademik telah banyak diteliti sebagai solusi untuk pemalsuan dokumen. Penelitian sebelumnya menunjukan penerapan privat blockchain mampu meningkatkan keamanan, efisiensi, dan integritas dokumen melalui penyimpanan nilai hash sebagai representasi keaslian data\footcite{KumarK2020}, sementara integrasi \textit{Decentralized Application}  (DApp), IPFS, dan \textit{QR Code} juga terbukti mendukung proses verifikasi dokumen yang lebih efisien\footcite{Gangwar2024}. Meskipun
		demikian, sebagian penelitian masih berfokus pada penyimpanan hash dan integrasi sistem, sehingga belum menerapkan data terstruktur seperti standar \textit{EIP-712}.


		\subsection{Blockchain dan \textit{Hyperledger Besu}}
		Blockchain merupakan teknologi buku besar terdistribusi (\textit{distributed ledger technology}) yang mencatat transaksi secara kronologis, permanen dan aman melalui mekanisme konsensus kriptografi\footcite{Zheng2017}. Dalam pengembangan sistem ini, yang digunakan adalah  \textit{Hyperledger Besu}\footcite{HyperledgerFoundation}, yaitu sebuah Ethereum open-source yang dirancang untuk mendukung kebutuhan pengguna pada tingkat perusahaan.  \textit{Hyperledger Besu} mendukung jaringan blockchain privat  yang memmungkinkan kendali penuh atas biaya transaksi (\textit{gas-fee}) serta penggunaan algoritma konsensus yang efisien seperti QBFT Hyperledger Foundation\footcite{HyperledgerFoundation}. Pada penelitian ini, jaringan privat dikonfigurasi dengan flag \textit{--min-gas-price=0} pada node validator. Sehingga transaksi dapat diproses tanpa membebankan biaya (\textit{gas}) kepada pengguna. 

		
		\subsection{Protokol EIP-712}
		EIP-712 (Ethereum Improvement Proposal) merupakan standar untuk hashing dan penandatanganan data terstruktur guna memberikan keamanan yang lebih baik untuk interaksi ke dompet kripto\footcite{Bloemen2017}. Standar ini memastikan pengguna melihat data yang bisa dibaca manusia saat memberikan tanda tangan digital, sehingga mencegah kesalahan manusia saat proses otorisasi data mentah yang tidak dapat terbaca.

		Prosedur penandatanganan dalam protokol ini mengikuti rumus pada Persamaan~\eqref{eq:eip712-signature}.
		\begin{equation}
			Signature = sign(keccak256(0x1901 \parallel domainSeparator \parallel hashStruct(message)))
			\label{eq:eip712-signature}
		\end{equation}
		Keterangan komponen rumus tersebut adalah:

		\begin{subsecenumerate}
			\item \textit{0x1901}: Merupakan prefix wajib yang terdiri dari byte kontrol 0x19 (untuk mencegah collision dengan transaksi RLP) dan version byte 0x01 yang spesifik untuk standar EIP-712.
			\item \textit{domainSeparator}: Nilai unik yang dihasilkan dari informasi aplikasi (nama, versi, chain ID, dan alamat kontrak). Hal ini menjamin bahwa tanda tangan untuk dokumen akademik di UNIDA Gontor tidak dapat digunakan ulang (replay attack) pada aplikasi atau jaringan blockchain lain.
			\item \textit{hashStruct(message)}: Representasi hash dari isi dokumen itu sendiri yang akan diverifikasi integritasnya.
		\end{subsecenumerate}

		\newpage
		Integritas pada level data struktur dikelola melalui fungsi hashStruct yang didefinisikan pada persamaan~\eqref{eq:eip712-hashstruct}.
		\begin{equation}
			hashStruct(s) = keccak256(typeHash \parallel encodedData(s))
			\label{eq:eip712-hashstruct}
		\end{equation}
		Penjelasan variable pada fungsi hashStruct:

		\begin{subsecenumerate}
			\item \textit{typeHash}: Merupakan hasil hashing dari definisi skema atau struktur data (misal: nama struct dan tipe data variabel di dalamnya). Hal ini memastikan bahwa penanda tangan menyetujui struktur data yang konsisten.
			\item \textit{encodedData(s)}: Merupakan hasil pengkodean nilai riil dari dokumen yang dikirim menjadi format 32-byte yang terstandarisasi.
		\end{subsecenumerate}

		
		\subsection{Integritas Data}
		Integritas data adalah jaminan bahwa data tetap akurat, konsisten, dan tidak mengalami perubahan yang tidak sah selama siklus hidupnya. Dalam sistem verifikasi dokumen, integritas menjadi pilar utama untuk memastikan bahwa informasi yang diterima oleh verifikator sama persis dengan yang diterbitkan oleh institusi asal. Teknologi blockchain menjamin integritas ini melalui penggunaan fungsi hash satu arah, dimana perubahan sekecil apapun akan menghasilkan nilai hash yang berbeda pada dokumen, sehingga modifikasi dapat terdeteksi secara instan\footcite{Zheng2017}. Sifat \textit{immutability} pada blockchain memastikan bahwa setelah data diverifikasi dan dimasukan kedalam blok, maka data tersebut menjadi dari bagian dari buku besar secara permanen. Dengan demikian, integritas data dalam penelitian ini tidak hanya mencangkup keamanan terhadap serangan luar, tetapi juga pencegahan terhadap manipulasi data internal oleh pihak yang memiliki akses.

		Verifikasi integritas berbasis blockchain dengan memanfaatkan fungsi hash kriptografi untuk menghasilkan sidik digital (\textit{digital fingerprint}) dari suatu berkas, kemudian menyimpan bukti hash tersebut pada blockchain sebagai referensi yang bersifat immutable. Pada saat proses verifikasi, sistem menghitung ulang nilai hash pada berkas kemudian membandingkannya dengan nilai hash yang telah tercatat pada blockchain\footcite{Ramya2026}. Pendekatan ini menjadi landasan dalam penelitian ini, dimana integritas data direpresentasikan melalui \textit{identityHash} dan integritas berkas PDF direpresentasikan sebagai \textit{fileHash}, yang kemudian dicocokan dan diverifikasi di jaringan blockchain.


		
		\subsection{Manajemen Kunci (\textit{Key Managment})}
		Manajemen kunci pada teknologi blockchain mencangkup seluruh siklus hidup kunci, mulai dari proses pembuatan hingga penyimpanan menggunakan keystore terenkripsi\footcite{Pal2021}. Implementasi ini krusial untuk menjamin fungsionalitas sistem yang tangguh dalam melindungi aset digital pengguna melalui mekanisme distribusi otoritas terukur guna mencegah eskalasi pihak luar yang tidak sah.


		\subsection{Solidity dan Smart Contract}
		Solidity adalah bahasa pemrograman tingkat tinggi berorientasi object yang dirancang khusus untuk pengembangan \textit{smart contract} pada platform ke blockchain yang kompatible dengan \textit{Ethereum Virtual Machine}  (EVM)\footcite{Wood2014}. Penggunaan \textit{Solidity} versi \textit{0.1.19} dalam penelitian ini menjamin stabilitas dan fungsionalitas yang luas terhadap protokol \textit{EVM} pada \textit{Hyperledger Besu}, terutama menghindari penggunaan instruksi (\textit{opcode}) yang belum didukung oleh konfigurasi jaringan privat blockchain.



		
	\section{Penelitian Terdahulu}
		\subsection{Manajemen Kunci dalam Sistem Blockchain}
		Manajemen kunci kriptografi merupakan aspek yang sangat penting dalam sistem berbasis blockchain, terutama dalam menjaga keamanan akses dan integritas data. Penelitian Om Pal dkk\footcite{Pal2021}. Menekankan pentingnya pengelolaan siklus hidup kunci (\textit{Key Lifecycle Management}) serta menggunakan skema \textit{Group Key Management} (GKM) untuk mendukung komunikasi yang aman dalam jaringan terdistribusi. Hasil penelitian ini menunjukan, bahwa pengelolaan kunci yang tersuktur mampu meningkatkan aspek kerahasiaan dan integritas data, serta mencegah akses yang tidak sah.

		Selain itu, penelitian dari Rafel Genes Duran dkk\footcite{Duran2020}. Mengusulkan pendekatan arsitektur \textit{middleware} yang memisahkan antra peran kunci antara otorisasi data dan pembayaran biaya transaksi (\textit{gas-fee}). Pendekatan ini mampu meningkatkan keamanan sekaligus memberikan kemudahan bagi pengguna dalam pengelolaan kunci kriptografi.

		Dalam penelitian ini, konsep manajemen kunci diadopsi melalui pemisahan peran antara \textit{signer} dan \textit{relayer} pada \textit{API Middleware}. Implementasi yang dilakukan masih bersifat sederhana, yaitu menggunakan penyimpanan berbasis \textit{environment variables} dan memori sistem (\textit{storage}) untuk pengelolaan kunci privat dimana kunci tersebut diamankan melalui proses enkripsi sebelum disimpan dan digunakan kembali pada proses penandatangan dokumen.

		
		\subsection{Sistem Verifikasi Dokumen Akademik Berbasis Blockchain}
		Pemanfaatan blockchain dalam verifikasi dokumen akademik telah banyak diteliti sebagai solusi terhadap pemalsuan dokumen. Penelitian oleh Kumar, dkk\footcite{KumarK2020}. Mengimplementasikan sistem verifikasi sertifikat berbasis \textit{private blockchain} berbasis Ethereum, dimana hash dokumen digunakan sebagai representasi keaslian data. Hasil penelitian menunjukan pendekatan ini lebih efisien dalam hal waktu dan biaya dibandingkan sistem berbasis blockchain publik, serta meningkatkan keamanan terhadap manipulasi data.

		Sementara itu, penelitian dari Gangwar dan Chaurasia\footcite{Gangwar2024} mengembangkan sistem autentikasi ijazah berbasis \textit{Decentralized Application} (DApp) dengan integrasi \textit{IPFS} sebagai media penyimpanan dan \textit{QR Code} untuk verifikasi fisik. Hasilnya menunjukan peningkatan signifikan dalam hal keamanan dan efisiensi biaya pembuatan dokumen.

		Penelitian-penelitian ini menunjukan bahwa blockchain efektif dalam menjamin keaslian sebuah dokumen akademik. Namun, sebagian besar masih berfokus pada penyimpanan hash atau integrasi sistem, sehingga membuka peluang untuk penelitian atau pengembangan lebih lanjut pada aspek penandatangan data terstruktur seperti \textit{EIP-712} yang menjadi fokus penelitian ini.

			
		\subsection{Keamanan dan Arsitektur Sistem Verifikasi Berbasis Blockchain}
		Aspek keamanan dalam sistem verifikasi dokumen akademik berbasis blockchain tidak hanya mencangkup integritas data, melainkan juga mencangkup autentikasi, otorisasi, privasi, dan kepemilikan data. Penelitian Saleh O. S. dkk\footcite{Saleh2020}. Mengembangkan kerangka kerja berbasis \textit{Hyperledger Fabric} yang mempunyai lima pilar utama, yaitu autentikasi, otorisasi, kerahasiaan, privasi, dan kepemilikan. Hasil penelitian menunjukan bahwa sistem yang dibangun mampu memastikan bahwa hanya pihak yang berwenang yang dapat menerbitkan dokumen, sementara pihak eksternal tetap dapat melakukan verifikasi secara transparan.
		
		Selain itu, Ramya dkk.\footcite{Ramya2026} mengusulkan mekanisme verifikasi integritas file dengan memanfaatkan fungsi hash kriptografi dan blockchain sebagai media penyimpanan nilai hash. Dalam penelitian tersebut, hash dari suatu berkas dihitung kembali dan dicatat pada blockchain, kemudian pada saat verifikasi akan dihitung kembali nilai hashnya dan membandingkannya dengan hash yang terimpan di blockchain

		Penelitian ini menjadi landasan dalam perancangan sistem yang tidak hanya berfokus pada integritas data, tetapi juga mempertimbangkan aspek keamanan secara menyeluruh. Dalam penelitian ini, pendekatan tersebut diadaptasi melalui penggunaan blockchain privat menggunakan \textit{Hyperledger Besu} dengan mekanisme kontrol akses dan validasi yang lebih terstruktur.

		% TABLE 
		\begin{landscape}

			\renewcommand{\arraystretch}{1.3}

			\begin{longtable}{|
				>{\centering\arraybackslash}p{0.8cm}|
				>{\centering\arraybackslash}p{3.5cm}|
				>{\centering\arraybackslash}p{1.2cm}|
				p{6cm}|
				p{5cm}|
				p{5cm}|}

				\caption{Perbandingan Penelitian Terdahulu}
				\label{tab:perbandingan_penelitian} \\

				\hline
				\textbf{No} &
				\textbf{Peneliti} &
				\textbf{Tahun} &
				\textbf{Judul Penelitian} &
				\textbf{Persamaan} &
				\textbf{Perbedaan} \\ \hline
				\endfirsthead

				\hline
				\endfoot

				\hline
				\endlastfoot

				1 &
				Rafael Genes Duran, dkk. &
				2020 &
				\textit{An Architecture for Easy Onboarding and Key Life-Cycle Management in Blockchain Applications} &
				Menerapkan pemisahan peran kunci (\textit{decoupling signature}) untuk kemudahan penggunaan. &
				Penelitian terdahulu fokus pada aplikasi web umum, sedangkan penelitian ini berfokus pada data akademik terstruktur menggunakan \textit{EIP-712}. \\ \hline

				2 &
				Saleh, O. S., dkk. &
				2020 &
				\textit{Blockchain Based Framework for Educational Certificates Verification} &
				Menjamin keamanan dan integritas sertifikat akademik melalui mekanisme otentikasi dan otorisasi. &
				Penelitian terdahulu menggunakan \textit{Hyperledger Fabric}, sedangkan penelitian ini menggunakan \textit{Hyperledger Besu} dengan protokol \textit{EIP-712}. \\ \hline

				3 &
				Om Pal, dkk. &
				2021 &
				\textit{Key Management for Blockchain Technology} &
				Berfokus pada keamanan manajemen kunci (\textit{Key Management}) untuk menjaga kerahasiaan data. &
				Penelitian terdahulu bersifat teoritis, sedangkan penelitian ini mengimplementasikan pemisahan peran Signer dan Relayer secara teknis. \\ \hline

				4 &
				Kumar, dkk. &
				2022 &
				\textit{Academic Certificate Verification System using Blockchain} &
				Menggunakan blockchain privat untuk verifikasi dokumen akademik. &
				Penelitian terdahulu belum menerapkan protokol \textit{EIP-712}, sedangkan penelitian ini menggunakan \textit{human-readable signing}. \\ \hline

				5 &
				Gangwar, S. dan Chaurasia, A. &
				2024 &
				\textit{Blockchain-based Authentication and Verification System for Academic Certificate} &
				Membangun DApp verifikasi ijazah berbasis blockchain dengan integrasi QR Code. &
				Penelitian terdahulu menggunakan Ethereum 2.0 yang memerlukan biaya gas, sedangkan penelitian ini menggunakan \textit{Hyperledger Besu} dengan mekanisme zero-gas. \\ \hline

				6 &
				Ramya &
				2026 &
				\textit{Blockchain Based Academic Certification Verification System} &
				Melakukan verifikasi integritas melalui perhitungan ulang hash dan penyimpanan bukti hash pada blockchain. &
				Penelitian terdahulu memverifikasi integritas file secara umum dan tidak menggunakan protokol \textit{EIP-712} untuk data terstruktur. \\
			\end{longtable}
		\end{landscape}